\documentclass[12pt]{article}

\usepackage[a4paper, margin=1in]{geometry}
\usepackage{amsmath, amssymb}
\usepackage{setspace}
\usepackage{hyperref}

\setstretch{1.2}

\title{\textbf{Biocentric Reality Engine: \\ A Neural Framework for the Internal Construction of Reality}}
\author{Iyer Ramkumar Ramasubramanian \\ \textit{Independent Researcher}}
\date{\today}

\begin{document}

\maketitle

\begin{abstract}
Traditional physical models assume an objective external reality independent of the observer. This paper introduces the \textbf{Biocentric Reality Engine}, a framework based on the core claim that reality is generated by the neural system itself. By integrating Evolutionary Graph CNNs, spiking neurons, and Integrated Information Theory ($\Phi$), we demonstrate a multi-agent system where "Reality" is not sensed, but rather co-created as a stable internal state projection among participating agents.
\end{abstract}

\section{Introduction}
Advances in neurobiology and statistical mechanics suggest that the observer is central to the structure of the universe. Following the principles of biocentrism, this research introduces a system where space and time are emergent properties of biological logic. Unlike the \textit{Homo Deus Brain Model}, which reacts to environmental signals, the \textbf{Biocentric Reality Engine} constructs its own environmental manifold.

\section{Theoretical Core}

The model is built upon four architectural additions to standard neural frameworks:

\begin{itemize}
\item \textbf{Evolutionary Graph CNN}: Simulates the biological "life" and structural growth of the network.
\item \textbf{Spiking Neurons}: Incorporates temporal dynamics to model the biological flow of time.
\item \textbf{IIT $\Phi$ with Default Mode Network}: Measures the integrated information and consciousness levels of the internal state.
\item \textbf{Free Will as Action Selection}: Redefines agency as the internal selection of outputs within the self-generated reality.
\end{itemize}

\section{System Architecture}

The system architecture replaces traditional sensing with internal projection:

\subsection{Reality Projection Layer}
Generates the "Reality Tensor," which represents the collective internal state projection of the agents rather than an external world.

\subsection{Integrated Information Layer}
Calculates $\Phi$ (consciousness) to stabilize the integration of data across the Default Mode Network.

\subsection{Holographic Boundary}
Consistent with the Holographic Principle, the projected reality is constrained by the information density of the internal neural weights.

\section{Mathematical Formulation}

The system focuses on the \textbf{Reality Mean} ($R_m$), representing the average value of the generated reality tensor across all agents. Stability is reached when:

\[
\Delta R_m \rightarrow 0
\]

In this state, agents have reached a **shared reality consensus**. The lack of an external reference point means $R_m$ is the only measurable reality within the system.

\section{Inference and Observations}

Data from the Biocentric Reality Engine indicates:
\begin{itemize}
\item \textbf{Internal Construction}: Reality is not being sensed from an external source; it is being actively constructed by the agents.
\item \textbf{Convergence vs. Fragmentation}: If the system converges (Reality Mean stabilizes), a stable world emerges. If it diverges, reality fragments into multiple subjective worlds.
\item \textbf{Biological Stability}: Similar to the HDBM, the system seeks a stable equilibrium, though here it is an equilibrium of perception rather than just action.
\end{itemize}

\section{Conclusion}
The Biocentric Reality Engine demonstrates that reality can be mathematically modeled as a measurement of a system's internal agreement. This shift from "sensing the world" to "constructing the world" aligns with a unified theory of biocentric intelligence.

\section{Repository for Experimentation}
\noindent
\url{https://github.com/spacetracker-collab/biocentrism_lanza}

\end{document}
